\chapter{Introduction}

In recent years, small satellites such as cube satellites (cubesats), and ride-share programmes has made research projects deployed in space viable for universities and other actors with limited resources. This has increased the availability of remote sensing data for earth observation. Amongst the available data is hyperspectral and multispectral imagery, which is also an important tool in areas such as medical research and food industry \cite{ientilucciAtmosphericCompensationHyperspectral2019}. Cubesats are often developed under strict budgets, and \acrshort{COTS} are used to reduce cost and development time. To facilitate the increased amounts of data, ongoing research is developing efficient solutions for processing and data acquisition for such missions.

As part of the \acrfull{hypso} project at the \acrfull{NTNU}, two 6U cubesats, HYPSO-1 and HYPSO-2, were launched in 2022 and 2024 respectively. These carry hyperspectral imagers capable of sensing wavelengths ranging from 430 nm to 800 nm over 120 spectral bands. Their primary purpose is to monitor algal blooms in the ocean, which can give indications as to the health of marine environments \cite{grotteOceanColorHyperspectral2022}. This is made possible by the reflection of characteristic wavelengths from chlorophyll, detectable by the on-board hyperspectral imagers. A nominal \acrshort{hypso} image has a size of roughly 150 MB. Images are combined with metadata such as georeferencing and anomaly detection, producing large amounts of data. This presents a challenge due to limited on-board resources for data storage, combined with low bandwidth of downlink over S-band radio. As an example, \acrshort{hypso}-1 is restricted to 75 MB of downlinked data per orbital revolution \cite{grotteOceanColorHyperspectral2022}. This prompts the need for on-board image compression to increase the scientific output and value of the mission.

Small satellites operate under strict limitations on on-board resources such as processor time and power budgets. This, combined with the complexity of image compression algorithms, makes the use of custom hardware accelerators ideal. For their implementation, \acrfull{FPGA}s are the perfect choice, providing great flexibility and speed at a low cost. The \acrshort{hypso} missions use a hardware accelerator of the CCSDS 123.0-B-1 compression algorithm which provides lossless compression. To provide even greater compression, the CCSDS 123.0-B-2 algorithm introduces near-lossless compression. A previous work by \citeauthor{vorhaugHyperspectralImageCompression2024} successfully implements a hardware accelerator for the new algorithm, but throughput is limited by computational data dependencies.

This thesis explores the use of a novel weight update scheme proposed by \citeauthor{jiaRemovalFeedbackLoop2025} for the CCSDS 123.0-B-2 compression algorithm. The purpose is to increase the throughput of the accelerator implemented by \cite{vorhaugHyperspectralImageCompression2024} for use on-board the \acrshort{hypso} missions. This should not come at significant degradation in compression performance or hardware resource utilization. Additionally, an open-source decompressor is implemented in order to verify the results.

The organization of this document is as follows. An introduction to the compression algorithm is given in \autoref{sec:background} before attempted implementation strategies of the open literature are presented in \autoref{sec:state-of-the-art}. \autoref{sec:previous-work} describes the work of \citeauthor{vorhaugHyperspectralImageCompression2024}. Suggested modifications, implementation and verification strategy is presented in \autoref{sec:implementation}. Finally, results are presented and discussed in sections \ref{sec:results} and \ref{sec:discussion}.
